%%%%%%%%%%%%%%%%%%%%%%%%%%%%%%%%%%%%%%%%%%%%%%%%%%%
\name{彭俊辉}
\basicContactInfo{18230688308}{pengjh0111@hnu.edu.cn}
\begin{tikzpicture}[remember picture, overlay]
  \node[anchor = north east] at ($(current page.north east)+(-1cm,-0.5cm)$)
    {\includegraphics[height=2.4cm]{avatar2}};
\end{tikzpicture}%
\vspace{-0.5cm}
%%%%%%%%%%%%%%%%%%%%%%%%%%%%%%%%%%%%%%%%%%%%%%%%%%%


%%%%%%%%%%%%%%%%%%%%%%%%%%%%%%%%%%%%%%%%%%%%%%%%%%%
\logosection{\faGraduationCap}{教育经历}
 
\datedline{%
  \textbf{湖南大学（985）} \quad 硕士 \quad 计算机技术 \quad
  主修课程：高等算法设计与分析 \key{91}，形式语言与自动机 \key{90}%
}{\dateRange{2024.09}{至今}}
 
\begin{itemize}
  \item 导师：\key{赵捷} \quad 研究方向：\key{AI编译器}；\key{算子编译与调度优化} \quad GPA：\key{3.55/4.00}
\end{itemize}
 
\datedline{%
  \textbf{湖南工商大学} \quad 本科 \quad 软件工程 \quad
  主修课程：计算机组成原理与汇编语言 \key{94}，算法设计与分析 \key{92}%
}{\dateRange{2020.09}{2024.06}}
 
\vspace{2pt}
%%%%%%%%%%%%%%%%%%%%%%%%%%%%%%%%%%%%%%%%%%%%%%%%%%%


%%%%%%%%%%%%%%%%%%%%%%%%%%%%%%%%%%%%%%%%%%%%%%%%%%%
\logosection{\faBriefcase}{实习经历}

\datedline{%
  \textbf{华为 2012 实验室（编译器实验室）} \quad 编译器开发实习生%
}{\dateRange{2026.06}{至今}}

\begin{itemize}
  \item 参与 \key{Triton-Ascend/AscendNPU-IR} 编译基础设施建设，聚焦 \key{Ascend NPU} 访存 IR 及端到端 Lowering 链路
  \begin{itemize}[leftmargin=1.8em, label=$\diamond$, topsep=1pt, itemsep=1pt]
    \item 重构 \key{hfusion.load/store} 及 Triton IR 至 AscendNPU-IR 的下降流程，以 hFusion Core Dialect 统一访存语义；将 MemRef/Bufferization 推迟至最终物化阶段，使中间层始终保持 \key{Tensor-level IR}
    \item 完善复杂访存模式支持：Store 覆盖 no-mask、masked、2D masked 与 boundary-check，Load 支持 Padding 下的边界访问，并补齐 Mask/Padding 在 AscendNPU-IR 中的表达与转换
  \end{itemize}
  \item 参与「面向昇腾的编译优化及 Agent 技术合作项目」，负责统一 \key{IR 抽象}方向的方案设计
\end{itemize}
%%%%%%%%%%%%%%%%%%%%%%%%%%%%%%%%%%%%%%%%%%%%%%%%%%%


%%%%%%%%%%%%%%%%%%%%%%%%%%%%%%%%%%%%%%%%%%%%%%%%%%%
\logosection{\faProjectDiagram}{科研与项目经历}

%% ── 项目一 ──────────────────────────────────────
\projecttitle{面向 STNN 的 Temporal-Tile 编译器 Chronos}{%
  投稿至 \key{ACM SIGOPS ATC 2026}，第一作者}

基于 \key{MLIR/ONNX-MLIR} 设计 STNN 并行调度系统，以 \key{tTile（Temporal Tile）} 抽象时序依赖，在兼容 \key{cuDNN/cuBLAS} 的同时挖掘 wavefront 并行性。

\begin{itemize}
  \item 设计 spatial/temporal 类型标注与 tTile 构建算法；经 \key{Loop Fusion/Coalescing/Tiling} 生成 CUDA Kernel，空间 Batch 算子 Lower 至 \key{cuDNN/cuBLAS/cuTENSOR}，保持 Vendor Library 兼容性
  \item 以 \key{4D 算子坐标}组织 inter-tTile wavefront 与 intra-tTile 多路径并行；在 11 个 STNN workload 上，相比 TVM、BladeDISC、Welder、TensorRT 实现 \key{1.25--2.04$\times$ 加速}，调度开销降低 \key{2--5 个数量级}
  \item 通过解析式评分在 $O(\log T)$ 个候选中选取最优 batch size，无需 Autotuning；以 \key{Track $\to$ CUDA Stream} 完成并行映射，并通过 CUDA Event 编排跨 Track、跨层级依赖
\end{itemize}
%% ── 项目二 ──────────────────────────────────────
\projecttitle{基于 MLIR 的 SNN 编译框架设计与实现}

基于 \key{MLIR} 构建 SNN 多层编译框架与端到端 Lowering Pipeline，支持模型从 ONNX 表示下降至目标芯片代码。

\begin{itemize}
  \item 扩展 \key{ONNX Dialect} 表示 SNN 算子，设计对接芯片 intrinsic 的 \key{Spike Dialect} 及 ONNX $\to$ Spike Lowering Pass
  \item 实现算子融合、循环变换等优化，并针对 \key{LIF/脉冲卷积} 完成 intrinsic 级算子优化与 \key{C + intrinsic} 代码生成
  \item 构建 \key{Python 前端}与自动代码生成链路，贯通模型导入、IR 转换、编译优化及目标代码生成
\end{itemize}
%%%%%%%%%%%%%%%%%%%%%%%%%%%%%%%%%%%%%%%%%%%%%%%%%%%


%%%%%%%%%%%%%%%%%%%%%%%%%%%%%%%%%%%%%%%%%%%%%%%%%%%
\logosection{\faCogs}{专业技能}
\begin{itemize}[parsep=0.5ex]
  \item \textbf{精通} \key{MLIR/LLVM 编译框架}，可独立设计 Dialect、开发优化 Pass 并构建多层 IR 转换 Pipeline；熟悉 GPU 后端优化：CUDA、Stream 并行、SM 调度、nsys / ncu profiling
  \item 熟悉主流编译优化技术：\key{Tiling}、\key{Loop Coalescing}、\key{Operator Fusion}、\key{Buffer Reuse}、\key{CSE}
  \item 具备 \key{ML编译与推理系统}（ONNX-MLIR、TVM、TensorRT）实践经验，支持模型从前端到后端部署流程
  \item 扎实的 \key{C/C++} 编程能力与数据结构基础；熟悉 \key{Linux / Git / Docker / CMake} 开发环境
\end{itemize}
%%%%%%%%%%%%%%%%%%%%%%%%%%%%%%%%%%%%%%%%%%%%%%%%%%%


\logosection{\faTrophy}{荣誉奖励}
\noindent
\key{PLMW@PLDI 奖学金}（PLDI 2025）$\;$\textbullet$\;$
\key{AsiaLLVM 学生奖学金}（LLVM Foundation）$\;$\textbullet$\;$
\key{硕士研究生二等学业奖学金}（湖南大学）$\;$\textbullet$\;$
\key{软件设计师中级证书}（国家工信部）$\;$\textbullet$\;$
\key{毕业生奖学金}$\;$\textbullet$\;$
\key{三好学生}（湖南工商大学）


%%%%%%%%%%%%%%%%%%%%%%%%%%%%%%%%%%%%%%%%%%%%%%%%%%%
\logosection{\faUniversity}{论文成果}

\noindent [1]\ Ke Yang, Guanghui Song, \textbf{Junhui Peng}, et al.
\textit{Efficient GPU Inference for Discrete Neural Networks}.
\key{MICRO 2026}

\smallskip
\noindent [2]\ \textbf{Junhui Peng}, et al.
\textit{Chronos: A Fast Temporal-Tile Scheduler for Spatiotemporal Neural Networks}.
第一作者，投稿至 \key{ACM SIGOPS ATC 2026}

\smallskip
\noindent [3]\ Shihan Yuan, Zuoyan Zhang, Guanghui Song, \textbf{Junhui Peng}, et al.
\textit{A Decoupled Analytical Model for Tile Size Selection in Affine Programs}.
\key{ACM TACO}
%%%%%%%%%%%%%%%%%%%%%%%%%%%%%%%%%%%%%%%%%%%%%%%%%%%
